\Needspace{14\baselineskip}
\CourseSection{Functional Description}\label{sec:functional}
% TODO (handoffs from Section 2 to other sections):
% - Requirements: no Section 3 requirement exists yet for the e-stop interface
%   (I15-I17), the GPS interface (I1), or the UWB tag update rate. Both are
%   marked with \reqTODO in tab:interfaces.
% - Requirements: MC-04 ("shall not energize its drive motors while the
%   emergency-stop control is engaged") conflicts with the Category 1 stop in
%   Section 4.2, where the motors stay energized to brake until the relay opens.
%   Suggested rewording in Section 6's handoff list.
% - Requirements: the trace tables allocate MC-02/04/05 to P-01, but the e-stop
%   is built in P-03.
% - Resources and Introduction: still list the existing golf push cart, motors,
%   motor controller, and battery. This section follows the Approach (hoverboard,
%   chassis, UWB, relay, e-stop receiver).
% TODO (open questions): confirm the hoverboard wheel diameter and mainboard type;
% select the e-stop radio; confirm whether the existing push-cart frame becomes
% the chassis; choose two-anchor ranging or AoA at TRK-2.

\CourseSubsection{System Overview and Design Alternatives}\label{sec:func-overview}

The \ac{ASCC} converts stored battery energy into motion toward one tagged golfer and decides, from radio, acoustic, and optical returns, whether that motion may continue. Club tracking and the distance display share the cart's power and processor but sit outside its motion path. The system boundary in Figure~\ref{fig:ascc-system-overview} includes the golfer's \ac{UWB} tag and the handheld e-stop. The golfer, safety operator, clubs, obstacles, satellites, and terrain are external. The main decisions are a used hoverboard with reflashed firmware as the drive base, a \ac{UWB} tag for tracking the golfer, a stop path that bypasses the Pi, ultrasonic sensors plus \ac{LiDAR} for obstacles, passive \ac{RFID} tags for clubs, and \ac{GPS} on the cart for distance. Table~\ref{tab:design-decisions} (Appendix~\ref{app:functional}) justifies each one against its alternatives and the needs that drive it.

\CourseSubsection{Functional Decomposition}\label{sec:func-decomp}

Figure~\ref{fig:functional} decomposes the system into six first-level functions: F1 move, F2 locate the golfer, F3 stop safely, F4 track clubs, F5 show distance, and F6 supply power. Each box performs one conversion of energy or information, and the labelled flows only carry its inputs and outputs. The separation of F3.4 from F3.7 and F3.8 is the central result. F3.4 runs on the Pi and brakes through I18, whereas F3.7 and F3.8 act on the energy path directly, so a hung Pi loses the brake but not the power cut. Table~\ref{tab:interfaces} (Appendix~\ref{app:functional}) specifies the flows that cross a hardware boundary and traces them to the interface requirements.

\begin{figure}[!htbp]
\centering
% Fully decomposed functional diagram. Interface labels trace to tab:interfaces.
\begin{tikzpicture}[x=1in,y=1in]
\def\ca{0.43}\def\cb{1.54}\def\cc{2.65}\def\cd{3.76}\def\ce{4.87}\def\cf{5.98}
\def\ra{0}\def\rb{-0.60}\def\rc{-1.20}\def\rd{-1.80}\def\re{-2.40}\def\rf{-3.00}\def\rg{-3.60}
% Row 1: distance
\node[extbox] (sat) at (\ca,\ra) {\acs{GPS} satellites};
\node[fnbox] (f51) at (\cb,\ra) {\textbf{F5.1} Compute position fix};
\node[fnbox] (f52) at (\cc,\ra) {\textbf{F5.2} Compute green distance};
\node[fnbox] (f53) at (\cd,\ra) {\textbf{F5.3} Compose screen};
\node[fnbox] (f54) at (\ce,\ra) {\textbf{F5.4} Convert frame to light};
\node[extbox] (read) at (\cf,\ra) {Golfer reads display};
% Row 2: clubs
\node[extbox] (club) at (\ca,\rb) {Clubs with tags};
\node[fnbox] (f41) at (\cb,\rb) {\textbf{F4.1} Read tag \acsp{UID}};
\node[fnbox] (f42) at (\cc,\rb) {\textbf{F4.2} Update inventory};
\node[fnbox] (f43) at (\cd,\rb) {\textbf{F4.3} Decide alert, club-out};
% Row 3: tracking
\node[extbox] (glf) at (\ca,\rc) {Golfer carrying tag};
\node[fnbox] (f21) at (\cb,\rc) {\textbf{F2.1} Emit \acs{UWB} pulses};
\node[fnbox] (f22) at (\cc,\rc) {\textbf{F2.2} Time pulses to ranges};
\node[fnbox] (f23) at (\cd,\rc) {\textbf{F2.3} Solve range, bearing};
\node[fnbox] (f24) at (\ce,\rc) {\textbf{F2.4} Compute $v,\omega$; zero if stale};
% Row 4: obstacles
\node[extbox] (obs) at (\ca,\rd) {Obstacles};
\node[fnbox] (f31) at (\cb,\rd) {\textbf{F3.1} Measure distance};
\node[fnbox] (f32) at (\cc,\rd) {\textbf{F3.2} Validate readings};
\node[fnbox] (f33) at (\cd,\rd) {\textbf{F3.3} Compare to stop zone};
% Row 5: e-stop
\node[extbox] (op) at (\ca,\re) {Safety operator};
\node[fnbox] (f35) at (\cb,\re) {\textbf{F3.5} Encode press as radio};
\node[fnbox] (f36) at (\cc,\re) {\textbf{F3.6} Decode; watch link};
% Arbiter spans rows 3 to 5
\node[fnbox, minimum height=1.60in, fill=LightGray!40] (f34) at (\cf,\rd) {\textbf{F3.4} Gate $v,\omega$; command brake on any stop request};
% Row 6: control
\node[fnbox] (f62) at (\ca,\rf) {\textbf{F6.2} Step down 25.2\,V to 5\,V};
\node[font=\tiny, text width=0.72in, align=left, inner sep=0pt] (rails) at (\cb+0.08,\rf) {to Pi, sensors, reader, display};
\node[fnbox] (f37) at (\cc,\rf) {\textbf{F3.7} Delay stop signal};
\node[fnbox] (f11) at (\cd,\rf) {\textbf{F1.1} Compute duty (\acs{FOC})};
\node[fnbox] (f14) at (\ce,\rf) {\textbf{F1.4} Sense rotor angle};
% Row 7: energy
\node[fnbox] (f61) at (\ca,\rg) {\textbf{F6.1} Store energy: 25.2\,V};
\node[fnbox] (f38) at (\cc,\rg) {\textbf{F3.8} Cut 25.2\,V: relay, kill switch};
\node[fnbox] (f12) at (\cd,\rg) {\textbf{F1.2} Invert \acs{DC} to 3-phase};
\node[fnbox] (f13) at (\ce,\rg) {\textbf{F1.3} Make torque: hub motors};
\node[extbox] (cart) at (\cf,\rg) {Cart and bag on terrain};
% Row 1 flows
\draw[sigarr] (sat) -- node[slab,above]{S1} (f51);
\draw[infarr] (f51) -- node[iflab,above]{I1} (f52);
\draw[infarr] (f52) -- node[iflab,above]{I2} (f53);
\draw[infarr] (f53) -- node[iflab,above]{I3} (f54);
\draw[sigarr] (f54) -- node[slab,above]{S2} (read);
% Position to club logic
\draw[infarr] ($(f51.south)+(0.25,0)$) -- ($(f51.south)+(0.25,-0.10)$) -- node[iflab,above,pos=0.5]{I4} ($(f43.north)+(-0.25,0.10)$) -- ($(f43.north)+(-0.25,0)$);
% Row 2 flows
\draw[sigarr] (club) -- node[slab,above]{S3} (f41);
\draw[infarr] (f41) -- node[iflab,above]{I5} (f42);
\draw[infarr] (f42) -- node[iflab,above]{I6} (f43);
\draw[infarr] (f43) -- node[iflab,right]{I7} (f53);
\draw[infarr] (f43.east) -- node[iflab,above]{I8} (\cf,\rb) -- (f34.north);
% Row 3 flows
\draw[sigarr] (glf) -- node[slab,above]{S4} (f21);
\draw[sigarr] (f21) -- node[slab,above]{S5} (f22);
\draw[infarr] (f22) -- node[iflab,above]{I9} (f23);
\draw[infarr] (f23) -- node[iflab,above]{I10} (f24);
\draw[infarr] (f24.east) -- node[iflab,above]{I11} (f24.east -| f34.west);
% Row 4 flows
\draw[sigarr] (obs) -- node[slab,above]{S6} (f31);
\draw[infarr] (f31) -- node[iflab,above]{I12} (f32);
\draw[infarr] (f32) -- node[iflab,above]{I13} (f33);
\draw[proarr] (f33.east) -- node[iflab,above]{I14} (f33.east -| f34.west);
% Row 5 flows
\draw[sigarr] (op) -- node[slab,above]{S7} (f35);
\draw[sigarr] (f35) -- node[slab,above]{S8} (f36);
\draw[proarr] (f36.east) -- node[iflab,above]{I15} (f36.east -| f34.west);
\draw[proarr] (f36) -- node[iflab,right]{I16} (f37);
\draw[proarr] (f37) -- node[iflab,right]{I17} (f38);
% Command to the drive
\draw[infarr] (f34.south) -- (\cf,-2.70) -- node[iflab,above,pos=0.3]{I18} (\cd,-2.70) -- (f11.north);
\draw[infarr] (f11) -- node[iflab,left]{I19} (f12);
\draw[infarr] (f14) -- node[iflab,above]{I20} (f11);
\draw[sigarr] (f13) -- node[slab,right]{S9} (f14);
% Energy chain
\draw[enarr] (f61) -- node[elab,above,pos=0.35]{E1} (f38);
\draw[enarr] (f38) -- node[elab,above]{E2} (f12);
\draw[enarr] (f12) -- node[elab,above]{E3} (f13);
\draw[enarr] (f13) -- node[elab,above]{E4} (cart);
\draw[enarr] (f61) -- node[elab,right]{E5} (f62);
\draw[enarr] (f62) -- node[elab,above]{E6} (rails);
% Legend
\begin{scope}[shift={(0,-4.08)}]
  \draw[enarr] (0,0) -- (0.32,0); \node[font=\tiny,anchor=west] at (0.34,0) {energy};
  \draw[infarr] (0.92,0) -- (1.24,0); \node[font=\tiny,anchor=west] at (1.26,0) {information};
  \draw[proarr] (2.03,0) -- (2.35,0); \node[font=\tiny,anchor=west] at (2.37,0) {protection};
  \draw[sigarr] (3.05,0) -- (3.37,0); \node[font=\tiny,anchor=west] at (3.39,0) {physical signal};
  \node[extbox, text width=0.22in, minimum height=0.13in] at (4.42,0) {};
  \node[font=\tiny,anchor=west] at (4.59,0) {outside the system};
\end{scope}
\end{tikzpicture}

\caption{Functional decomposition. Labels E (energy), I (information), and S (physical signal) identify interfaces; Table~\ref{tab:interfaces} (Appendix~\ref{app:functional}) traces them to requirements.}
\label{fig:functional}
\end{figure}

\CourseSubsection{Logical Decomposition and Interfaces}\label{sec:func-logical}

Figure~\ref{fig:logical} (Appendix~\ref{app:functional}) allocates each function in Figure~\ref{fig:functional} to one logical element and states the action between elements. The following controller, safety controller, club inventory logic, and display composition are separate logical elements but share one physical processor, so a Pi fault is a common cause across them. This is why the safety interlock must act on the drive unit without passing through the safety controller, and why the mainboard enforces its own limits.

\CourseSubsection{Physical Architecture}\label{sec:func-physical}

Figure~\ref{fig:physical} (Appendix~\ref{app:functional}) places the components on the chassis. The hoverboard's two halves are locked at the pivot to form a rigid drive axle, with a front caster added. The pack and mainboard stay low to limit tip-over (H-05), and the obstacle sensors sit low at the front to see low objects (H-02). The \ac{UWB} anchors are as far apart as the cross-bar allows, because bearing error scales with $1/b$. The \ac{GPS} antenna is high on the handle and the sensitive antennas are kept away from the hub motors (C-4, C-5). Interfaces between the Pi and the rest of the cart use the buses in Table~\ref{tab:interfaces}, and the compatibility checks at INT-1 confirm they do not conflict.

\CourseSubsection{Prototype Definitions}\label{sec:func-prototypes}

The prototypes are P-01 drive base, P-02 tracking and following, P-03 safety layer, P-04 smart bag, P-05 distance display, and P-06 integrated \ac{ASCC}. Table~\ref{tab:prototype-ids} (Appendix~\ref{app:functional}) links each to its functions, physical elements, and tests, and Table~\ref{tab:prototypes} gives its development stages.
